
% Hilfsdefinition für linksbündige Spalten mit fester Breite und Umbruch
\newcolumntype{L}[1]{>{\raggedright\arraybackslash}p{#1}}

% Start der xltabular-Umgebung
% {\textwidth} = nutzt die volle Seitenbreite
% L{4.5cm} = erste Spalte fix 4.5cm, linksbündig
% X = mittlere Spalte nutzt den restlichen Platz flexibel
% L{3.5cm} = letzte Spalte fix 3.5cm, linksbündig
\begin{xltabular}{\textwidth}{L{4.5cm} X L{3.5cm}}
    
    \caption{Übersicht von verwandten Arbeiten, sortiert nach Erscheinungsjahr} \label{tab:literaturuebersicht} \\
    
    \toprule
    \textbf{Titel (Autoren)} & \textbf{Relevanz für die Arbeit} & \textbf{Geeignet für Kapitel} \\
    \midrule
    \endfirsthead % Kopfzeile nur auf der ersten Seite
    
    \multicolumn{3}{c}{\textit{Fortsetzung von der vorherigen Seite}} \\
    \toprule
    \textbf{Titel (Autoren)} & \textbf{Relevanz für die Arbeit} & \textbf{Geeignet für Kapitel} \\
    \midrule
    \endhead % Kopfzeile auf allen Folgeseiten
    
    \bottomrule
    \multicolumn{3}{r}{\textit{Weiter auf der nächsten Seite}} \\
    \endfoot % Fußzeile auf allen umgebrochenen Seiten
    
    \bottomrule
    \endlastfoot % Fußzeile nur am Ende der gesamten Tabelle

    % --- TABELLENINHALT START ---
    
    A Comprehensive Review on File Containers-Based Image and Video Forensics (Yang u.\,a. 2025)\cite{yang_comprehensive_2025} & 
    Aktueller Überblick über Container-Forensik & 
    Container-Forensik \\ \midrule
    
    Media source similarity hashing (MSSH): A practical method for large-scale media investigations (Klier und Baier 2025)\cite{klier_media_2025} & 
    Similarity Hashing für JPEG samt GitHub & 
    Similarity Hashing \\ \midrule
    
    Learning Hierarchical Fingerprints via Multi-Level Fusion for Video Integrity and Source Analysis (Y. Li u.\,a. 2024)\cite{li_learning_2024} & 
    Argumentiert, dass Container-Spuren durch Re-Muxing leicht manipulierbar sind & 
    Container-Forensik \\ \midrule
    
    Forensic Analysis of Binary Structures of Video Files (Hasan u.\,a. 2021)\cite{hasan_forensic_2021} & 
    Untersucht binäre Unterschiede in ftyp- und moov-Boxen zwischen Kameraherstellern & 
    Aufbau MP4, Container-Forensik \\ \midrule
    
    Forensic Analysis of Video Files Using Metadata (Xiang u.\,a. 2021)\cite{xiang_forensic_2021} & 
    Nutzt Vektorisierung von Metadaten-Bäumen für MP4 & 
    Aufbau MP4, Container-Forensik \\ \midrule
    
    Digital Video Source Identification Based on Container's Structure Analysis (Ramos López u.\,a. 2020)\cite{ramos_lopez_digital_2020} & 
    Nutzt hierarchische Container-Strukturen diverser Video-Dateien zur unüberwachten Quellenerkennung anhand der VISION- und ACID-Datensätze & 
    Aufbau MP4, Container-Forensik \\ \midrule
    
    Efficient Video Integrity Analysis Through Container Characterization (Yang u.\,a. 2020)\cite{yang_efficient_2020} & 
    Etabliert einen quantitativen Ansatz mittels Bag-of-Words-Modell und Entscheidungsbäumen zur Integritätsprüfung. Liefert zudem den EVA-7K-Datensatz (manipulierte und komprimierte Social-Media-Videos auf Basis „VISION“ von Shullani u.\,a.) & 
    Container-Forensik, Video-Datensätze \\ \midrule
    
    Information technology - Coding of audio-visual objects - Part 14: MP4 file format (ISO 2020)\cite{iso_14496-14_2020} & 
    ISO-Norm, definiert die formale Spezifikation und hierarchische Objektstruktur (Atome/Boxen) des MP4-Formats & 
    Aufbau MP4 \\ \midrule
    
    A Video Forensic Framework for the Unsupervised Analysis of MP4-Like File Container (Iuliani u.\,a. 2019)\cite{iuliani_video_2019} & 
    Formalisiert MP4 als Baumstruktur & 
    Aufbau MP4, Container-Forensik \\ \midrule
    
    VISION: a video and image dataset for source identification (Shullani u.\,a. 2017)\cite{shullani_vision_2017} & 
    Liefert Datensatz für native Videoquellen (Dataset „VISION“), Grundlage für Social-Media-Variante (Dataset „EVA-7K“ von Yang u.\,a.) & 
    Video-Datensätze \\ \midrule
    
    Forensic analysis of video file formats (Gloe u.\,a. 2014)\cite{gloe_forensic_2014} & 
    Zeigen auf, dass Hersteller und Software-Encoder spezifische Spuren in der Struktur von AVI- und MP4-Dateien hinterlassen, etwa durch die Anordnung und Benennung von Datenblöcken & 
    Aufbau MP4, Container-Forensik \\
    
    % --- TABELLENINHALT ENDE ---

\end{xltabular}